\documentclass[12pt]{article}
\usepackage{lineno}
%\usepackage{ametsoc}
\usepackage{natbib}
\usepackage{a4wide}
\usepackage{longtable}
\usepackage{amsfonts}
\usepackage{amssymb}

\usepackage{color}
\usepackage[pdftex]{graphicx}
\usepackage[pdftex,colorlinks, citecolor=black,linkcolor=black,urlcolor=black]{hyperref}
\definecolor{darkgreen}{rgb}{0.1,0.4,0.1}
\definecolor{darkblue}{rgb}{0.1,0.1,0.3}



\newif\ifdetail
\global\detailfalse
%\global\detailtrue

\def\bem#1{\colorbox{green}{\em #1}}

\newcommand{\beq}  { \begin{eqnarray} }
\newcommand{\eeq}  { \end{eqnarray}}
\newcommand{\beeq}  { \begin{eqnarray*} }
\newcommand{\eeeq}  { \end{eqnarray*}}
\newcommand{\eq }  [1]{Eq.~(\ref{#1})}  % Eq. (#1)
\newcommand{\eqs }  [1]{Eqs.~(\ref{#1})}
\newcommand{\sect}  [1]{Sec.~(\ref{#1})}  % Eq. (#1)
\newcommand{\fig}  [1]{Fig.~\ref{#1}}   % Fig. #1
\renewcommand{\v}[1]{{\mbox{\boldmath$ #1 $}}}  % Vektor durch Fettdruck
\newcommand{\m}    [1]{{\bf  #1 }}              % Matrix in sans serif
\newcommand{\gz}[1] {{{d #1} \over {dz}}}
\newcommand{\pt}[1]{{\partial_t #1}}
\newcommand{\ptt}[1]{{\partial_{tt} #1}}
\newcommand{\pttt}[1]{{\partial_{ttt} #1}}
\newcommand{\pttx}[1]{{{\partial^3 #1} \over {\partial t^2 \partial x}}}
\newcommand{\ptty}[1]{{{\partial^3 #1} \over {\partial t^2 \partial y}}}
\newcommand{\ptx}[1]{{{\partial^2 #1} \over {\partial t \partial x}}}
\newcommand{\pty}[1]{{{\partial^2 #1} \over {\partial t \partial y}}}
\newcommand{\pz}[1]{{\partial_z #1}}
\newcommand{\pzz}[1]{{\partial_{zz} #1}}
\newcommand{\px}[1]{{\partial_x #1}}
\newcommand{\py}[1]{{\partial_y #1}}
\newcommand{\pxi}[1]{{\partial_i #1}}
\newcommand{\pxj}[1]{{\partial_j #1}}
\newcommand{\pxk}[1]{{\partial_k #1}}
\newcommand{\pb}[1]{{{\partial #1} \over {\partial b}}}
%\newcommand{\pxy}[1]{{{\partial^2 #1} \over {\partial x \partial y}}}
\newcommand{\pxy}[1]{\partial_{xy} #1}
\newcommand{\pyy}[1]{\partial_{yy} #1}
\newcommand{\pxx}[1]{\partial_{xx} #1}
\newcommand{\kx}{{\bf k} \times}
\newcommand{\dt}[1]{{{d #1} \over {d t}}}
\newcommand{\Dt}[1]{{{D #1} \over {D t}}}
\newcommand{\Div}[1]{ \nabla \cdot #1 }
\renewcommand{\div}[1]{ \nabla_{yz} \cdot #1 }
\newcommand{\Bar}[1]{ \overline{ #1} }
\newcommand{\rvec} [1]{
    \raisebox{-1.5ex}{$\stackrel{\textstyle #1}{\neg}$} }% rotierter Vektor
\newcommand{\nablq}   {\rvec\nabla}                    % rotated nabla operator
\newcommand{\Vec}[1]{\left(\begin{array}{c} #1 \end{array}\right) }
%\renewcommand{\mho}{\mathcal{W}}
\newcommand{\pn}[1]{{{\partial #1} \over {\partial n}}}
\newcommand{\pnn}[1]{{{\partial #1} \over {\partial m}}}
\newcommand{\ps}[1]{{{\partial #1} \over {\partial s}}}
\newcommand{\p}{{\partial}}
\newcommand{\vn}{{\v \nabla}}
\newcommand{\lk}{\left(}
\newcommand{\rk}{\right)}

\newcommand{\tl}{\tilde}

\newcommand{\I}{\mathcal{I}}
\newcommand{\E}{\mathcal{E}}
\newcommand{\A}{\mathcal{A}}

\newcommand{\ep}{\epsilon}

\begin{document}

 \section*{The simple Primitive equation model}

Consider the Boussinesq equations in hydrostatic approximation
\beq \label{eq1}
  \pt{\v u } +   f \rvec{\v u}  +   \vn p =   - \v u \cdot \vn \v u - w \pz{} \v u +  \v F( \v u)~~,~~
    \pt{b}   =  - \v u \cdot \vn b  - w \pz{} b +  M b
     \eeq
  with the diagnostic relations     $\pz p =b$ and $\vn \cdot \v u + \pz w =0$.
  Vectors are two-dimensional and $\rvec{\v u} $ denotes  anticlockwise
 rotation by $90^o$. 
  Pressure is scaled with background density $\rho_0$ and thermodynamics have been
  simplified to a buoyancy equation.
  $\v F$  and $M$ are the friction and mixing operators.
   At the free surface $z=\eta$, we also have the
      condition 
      \beq \pt p_s  = -   g \vn \cdot  \int_{-H}^\eta \v u dz 
       \eeq
  with $p_s = p|_{z=\eta} =  g \eta$.
  Alternatively, we use a rigid lid with $w=0$ at $z=0$
  and $p_s = p|_{z=0}$.
We split full pressure $p$ into hydrostatic part and surface pressure.
\beq
p = p_h(z) + p_s ~,~
p_h =  \int_{-h}^z b  dz'
\eeq
where  $p_s$ is to be determined by the external mode solvers.



\subsection*{Simple implicit external solver for rigid lid}

Velocity and buoyancy and pressure are all co-located in time in this scheme.
The buoyancy equation is integrated as
\beq
  b^{n+1}- b^n = \Delta t \lk  \Bar{\delta b}^{n+1/2} + M (b^{n})    \rk 
\eeq
with the advective buoyancy tendency
\beq
\delta b^n =  - \vn \cdot \v u^{n} b^n - \p_z w^{n} b^n ~~,~~
\Bar{ \delta b}^{n+1/2} = A \delta b^n + B \delta b^{n-1} + C \delta b^{n-2}
 \eeq
 interpolated with the Adam-Bashforth (AB3) scheme  to time level $n+1/2$.
 The biharmonic horizontal mixing (attribute \verb|Khbi|) is treated with an Euler forward step.
 Vertical mixing follows the options described in the section on vertical mixing and friction below.
 
An intermediate velocity $\v u^*$ is calculated from
\beq
\v u^* -\v  u^n = \Delta t \lk    \Bar{\delta \v u}^{n+1/2}   + \v F (\v u^{n}) \rk 
\eeq  
with the momentum tendency $\delta \v u$ excluding the surface pressure gradient and friction
\beq
 \delta \v u^n =  - f \rvec{\v u}^n - \vn  p_h^n  - \v u^n \cdot \vn \v u^n - w^n \pz{} \v u^n 
\eeq
These tendencies are extrapolated with AB3  to time level $n+1/2$ 
\beq
\Bar{ \delta \v u}^{n+1/2} = A \delta \v u^n + B \delta \v u^{n-1} + C \delta \v u^{n-2}
\eeq
Friction is treated with an Euler forward step.
From the vertically integrated continuity equation follows the condition
\beq
\vn \cdot  \int_{-h}^0  \v u^*  =  \Delta t  \vn \cdot h \vn p_s^{n+1}
\eeq
which is solved with a conjugate-gradient iteration for $p_s^{n+1}$
(attributes \verb|epsilon| for the convergence threshold and
\verb|max_iterations|; a BiCGStab variant is implemented as well and is
selected by editing the solver call in \verb|pressure.py|).
The iteration count and the estimated error are available per step as
\verb|congr_itts| and \verb|estimated_error|. The intermediate velocity is  then
corrected as
\beq
 \v u^{n+1} = \v u^* - \Delta t \vn p_s^{n+1}
\eeq
It is possible to include an implicit free surface in the time stepping.
Vertical velocity is calculated from continuity
\beq
 w^{n+1} = -  \int_{-h}^z \vn \cdot \v u^{n+1}
\eeq
This scheme is also used as default in the MITgcm.

\subsection*{Simple split-explicit external mode solver}


   
Velocity is known in this scheme at $t=n+1/2$, while $b$ and $p$ at level $n$.
The momentum equation is discretised as
\beq
\v u^* -\v  u^{n-1/2} = \Delta t \lk - f  \Bar{\rvec{\v u}}^n+   \Bar{\delta \v u}^{n}  - \vn  p_h^n - \vn p_s^n+ \v F (\v u^{n-1/2}) \rk 
\eeq  
with the intermediate velocity $\v u^*$ and 
AB3 extrapolation of Coriolis force and advection terms tendencies
\beq
 \Bar{\rvec{\v u}}^n =  A \rvec{\v u}^{n-1/2} + B \rvec{\v u}^{n-3/2} + C \rvec{\v u}^{n-5/2}
~~,~~ \Bar{ \delta \v u}^{n} = A \delta \v u^{n-1/2} + B \delta \v u^{n-3/2} + C \delta \v u^{n-5/2}
\eeq
and
\beq
 \delta \v u^{n-1/2} =     - \v u^{n-1/2} \cdot \vn \v u^{n-1/2} - w^{n-1/2} \pz{} \v u^{n-1/2}
\eeq
For the hydrostatic and surface pressure $p_h+p_s$ no extrapolation is made, since they are at level $n$ anyways, 
and friction is treated with Euler forward.

Now the momentum tendencies are vertically integrated but without
surface pressure and Coriolis contributions.
\beq
 \v R^n = \int_{-h}^0  \lk  \Bar{\delta \v u}^{n}  - \vn  p_h^n+ \v F(\v u^{n-1/2}) \rk dz 
\eeq
The external mode is integrated with sub-cycling over $m=1, ...M$
($M$ set by the attribute \verb|sub_cycling|) from $n$ to $n+1$ with
a simple, dissipative  scheme proposed by Demange et al (2019) using the forcing $\v R^n$
\beq
\v U^{n+(m+1)/M} - \v U^{n+m/M} &=& \Delta t/M \lk - f  \rvec{\Bar{\v U}}^{n+(m+1/2)/M} -  h \vn p_s^{n+m/M} + 
 \v R^n  \rk
 \\
 p_s^{n+(m+1)/M} - p_s^{n+m/M} &=& - g \Delta t/M \lk  (1+\theta)  \vn \cdot \v U^{n+(m+1)/M} 
 - \theta   \vn \cdot \v U^{n+m/M}
 \rk
\eeq
with the transport $\v U = \int_{-h}^0 \v u\, dz$, and $\Bar{\v U}^{n+(m+1/2)/M}$ interpolated
with AB2.   Demange et al (2019)  found  $\theta=0.14$ to be the best choice in terms of phase error and damping.
In the model $\theta$ is the attribute \verb|fb_theta| with default $0.14$.
Larger values damp grid-scale barotropic waves more strongly.

The sub cycle starts from the transport $\v U^n$ and the surface pressure $p_s^n$ which the
sub cycle of the previous time step ended with.
This transport is model state (\verb|U_bt|, \verb|V_bt|) and part of the restart.
It is computed from the vertical integral of $\v u^n$ only at a cold start.
Two transports result from the sub cycle.
The transport at the end of the sub cycle, $\v U^{n+1} = \v U^{n+M/M}$, starts the next sub cycle.
The average over the sub cycle
\beq
\tl{\v U}^{n+1}=\frac{1}{M} \sum_{m=1}^{M} \v U^{n+m/M}  + \frac{\theta}{M} \lk  \v U^{n+1} - \v U^n  \rk
\eeq
is used to correct the barotropic part of the intermediate velocity
\beq
  \v u^{n+1/2} = \v u^* -1/h \lk  \int_{-h}^0 \v u^* dz - \tl{\v U}^{n+1} \rk
\eeq
The next sub cycle must not start from the average $\tl{\v U}^{n+1}$.
That would pair the averaged transport with the surface pressure at the end of the sub cycle.
This mismatch grows from step to step when the Courant number of the external
gravity wave is large (Banerjee et al.\ 2024, the split-explicit solver of FESOM2).
The buoyancy equation is integrated using $\v u^{n+1/2}$ and $w^{n+1/2} = - \int_{-h}^z \vn \cdot \v u^{n+1/2}$
as
\beq
  b^{n+1}- b^n = \Delta t \lk  \delta b^{n+1/2} + M (b^{n})    \rk 
\eeq
with the buoyancy tendency
\beq
\delta b^n =  - \vn \cdot \v u^{n+1/2} b^n - \p_z w^{n+1/2} b^n
 \eeq
 interpolated with AB3 to time level $n+1/2$.
 This scheme is one option in FESOM2.

In both schemes the advective tracer flux through the surface face is zero by default, as in pyOM2.
Under the linear free surface the vertical velocity $w$ at the surface face is not zero.
A tracer that is uniform in space then changes in the surface cell at the rate $ b\, w|_{z=0} / \Delta z = - b\, \vn \cdot \int_{-h}^0 \v u\, dz / \Delta z$.
With the attribute \verb|free_surface_tracer_flux| the flux through the surface face is
$w\, b$ instead, with $w$ at the surface face and $b$ of the surface cell.
A uniform tracer then stays uniform.
This option is off by default.
The external-mode treatment is selected by the attributes
\verb|implicit_free_surface| and \verb|explicit_free_surface|: both false
gives the rigid lid, the first the implicit free surface (both solved with
the iterative solver above), the second the split-explicit scheme.
  

\section*{Vertical mixing and friction options}

Vertical mixing of buoyancy and vertical friction are organised in a
hierarchy of options similar to pyOM2, selected by the attribute
\verb|vmix_scheme| $\in$ \{\verb|'explicit'|, \verb|'implicit'|\} and two
switches (module \verb|implicit_vmix.py|; the previous flag
\verb|enable_implicit_vmix| maps to the implicit scheme with convective
adjustment, so existing configurations are unchanged).

\emph{Explicit scheme.} Fixed coefficients $\kappa$ (buoyancy) and $A_v$
(momentum) only, integrated forward in time,
\beq
 b^{n+1} = b^n + \Delta t\, \pz{} \lk \kappa\, \pz b^n \rk, \qquad
 \v u^{n+1} = \v u^{n+1} + \Delta t\, \pz{} \lk A_v\, \pz{\v u} \rk
\eeq
subject to the diffusive stability limit, which is checked at start-up
(\verb|check_vmix_stability|, called from \verb|prepare_for_run|):
$ 4 \max(\kappa, A_v)\, \Delta t / \Delta z^2 \le 1/2$; the model refuses to
run otherwise, and also refuses convective adjustment or the TKE closure
under the explicit scheme.

\emph{Implicit scheme.} Both operators are integrated with backward Euler
(unconditionally stable, so arbitrarily large coefficients are allowed),
as a tridiagonal solve per column with the coefficient located on the
$W$ points. Three sub-options for the coefficients:
(i) constant $\kappa$ and $A_v$;
(ii) convective adjustment (\verb|enable_conv_adjust|): background
$\kappa$, but $\kappa = \kappa_{conv}$ (large) where $N^2 < 0$;
(iii) the TKE closure (\verb|enable_tke|), which overrides all constant
coefficients for both buoyancy and momentum as described below.
Vertical friction is applied after the surface pressure correction: since
the vertical operator with zero-flux boundary conditions leaves the
column-integrated momentum unchanged, the barotropic mode is not
perturbed, and $w$ is re-diagnosed afterwards.

\emph{Bottom friction.} Two options act on the velocity $\v u_b$ in the lowest wet cell
of each $U$ and $V$ column. Both are off by default.
Linear friction (\verb|enable_bottom_friction|) has the rate $r = r_{bot}$ in 1/s
(attribute \verb|r_bot|).
Quadratic friction has the rate
$r = C_d |\v u_b| / \Delta z$ with the drag coefficient $C_d$ (attribute \verb|r_quad_bot|).
It is switched on by
\verb|enable_quadratic_bottom_friction|.
Here $|\v u_b|$ is the speed in the bottom cell. The other velocity component is averaged
to the $U$ or $V$ point from the four surrounding points.
If both options are on, the two rates add.
The friction is an implicit step on the bottom cell,
\beeq
 \v u_b^{new} = \v u_b / \lk 1 + \Delta t\, r \rk
\eeeq
In the rigid-lid and implicit free-surface scheme it is applied after the vertical friction.
In the split-explicit scheme it is applied before the external mode is integrated.
There the same bottom stress also acts on the transport inside the sub cycle.
Every sub step divides $\v U$ by $1 + (\Delta t/M)\, r\, \Delta z / h$.

\subsection*{TKE closure}

The closure of Gaspar et al.\ (1990) is ported from pyOM2
(\verb|tke.py|, mixin class). All closure constants are configuration
attributes. They are set before \verb|init_tke|.
The defaults are the pyOM2 \verb|tke_module| defaults $c_k = 0.1$, $c_\epsilon = 0.7$, $\alpha = 1$,
$\ell_{min} = 10^{-12}\,$m and $\kappa_M^{max} = 100\,$m$^2$/s.
The floors $\kappa_M^{min}$ and $\kappa_H^{min}$ of viscosity and diffusivity are zero by default, as in pyOM2.
With the TKE closure on, the constant coefficients $A_v$ and $\kappa$ of the hierarchy above are not used.
A background viscosity or diffusivity is therefore a choice of the configuration.
It is made by setting the floors. Note that the
pyOM2 configurations of the acc family deviate from the module defaults and
set $\alpha = 30$ (a proxy for vertical TKE transport by pressure work,
which spreads the surface-injected TKE substantially deeper) and
$\ell_{min} = 10^{-8}\,$m; for consistency with those experiments both
should be set accordingly (\verb|alpha_tke|, \verb|mxl_min|). The turbulent
kinetic energy $e$ is a prognostic variable on the $W$ points,
\beq
 \pt e = \pz{} \lk \alpha\, \kappa_M \pz e \rk
 + K_{diss} - P_{diss} - c_\epsilon\, e^{3/2} / \ell
\eeq
with the mixing length $\ell = \sqrt{2 e}/N$, limited such that it cannot
grow faster than the distance to the boundaries (two sweeps, pyOM2's
\verb|tke_mxl_choice=2|), and the diffusivities
\beq
 \kappa_M = \min( \kappa_M^{max},\, c_k\, \ell \sqrt{e} ), \qquad
 \kappa_H = \kappa_M / \max(1, \min(10,\, 6.6\, Ri)), \qquad
 Ri = N^2 \kappa_M / K_{diss}
\eeq
There is no separate convective branch.
Where $N^2 < 0$ the Richardson number is negative.
The Prandtl number is then one and $\kappa_H = \kappa_M$, exactly as in pyOM2.
This replaces the convective-adjustment option.
The ratio $K_{diss}/\kappa_M$ in $Ri$ is the squared vertical shear of the friction operator.
With \verb|mix_eulerian| (see below) it is $\pz{\bar{\v u}} \cdot \pz{} \lk \bar{\v u} + \v u_S \rk$.
Without it, it is the squared shear of the model velocity.
$N^2$ lives on the $W$ points, i.e.\ on the top face of each cell.
At the surface face there is no cell above, so $N^2$ is copied from the face below, as in pyOM2.
For \verb|eq_of_state| $=5$ the density of the upper cell is evaluated at the pressure of the
lower cell and compared with the density of the lower cell (adiabatic displacement, as in pyOM2).

Two sets of limits for the mixing length are available. They are selected by \verb|mxl_choice|.
The default \verb|mxl_choice| $=1$ is the pyOM2 form named above.
Its sweeps run in the order of \verb|tke.f90|.
First a downward sweep from the surface limits $\ell_k \le \ell_{k+1} + \Delta z$.
It starts from the unlimited value at the surface face.
Then the value at the surface face is limited to $\ell \le \ell_{min} + z_0 + \Delta z$,
with the roughness length $z_0$ given below.
Finally an upward sweep from the bottom limits $\ell_k \le \ell_{k-1} + \Delta z$.
The same length enters the diffusivity and the dissipation.
With \verb|mxl_choice| $=2$ the limits of NEMO 4.2 are used (\verb|nn_mxl| $=3$ with \verb|ln_mxl0|).
All lengths live on the $W$ points, i.e.\ at the top face of each tracer cell.
The index $n$ counts these points upward. Here $n=1$ is the lowest point of the array
and $n_s$ the surface face.
The starting point is the buoyancy length
\beeq
 \ell_n = \max \lk \sqrt{2 e} / \sqrt{\max(N^2, 10^{-12}\,\mbox{s}^{-2})},\ \ell_{min0} \rk
\eeeq
with $\ell_{min0} = 0.04\,$m (attribute \verb|mxl_min0|).
This floor is applied at all $W$ points, also at dry points where $e=0$.
The surface value is then reduced to the surface length, $\ell_{n_s} = \min(\ell_{n_s}, \ell_0)$, with $\ell_0$ given below.
Since $\ell_0$ is set before the downward sweep, it bounds the whole near-surface profile and not only the surface point.
Two profiles are formed from $\ell_n$. Going downward from the surface the length is limited by the distance to the surface,
\beeq
 \ell^{up}_n = \min \lk \ell_n,\ \ell^{up}_{n+1} + \Delta z \rk, \quad n = n_s-1, \ldots, 1, \qquad \ell^{up}_{n_s} = \ell_{n_s}
\eeeq
Going upward from the bottom of the array the length is limited by the distance to the bottom,
\beeq
 \ell^{dwn}_n = \min \lk \ell_n,\ \ell^{dwn}_{n-1} + \Delta z \rk, \quad n = 2, \ldots, n_s, \qquad \ell^{dwn}_1 = \ell_1
\eeeq
The two sweeps run over the whole array, dry points included.
In a column whose lowest wet cell lies above the bottom of the array, the $W$ point below the lowest wet cell is dry.
There $\ell_n = \ell_{min0}$, so the lowest wet point gets $\ell^{dwn} \le \ell_{min0} + \Delta z$.
In a column that is wet down to the bottom of the array the lowest point gets $\ell^{dwn}_1 = \ell_1$ without this limit.
The diffusivity uses the geometric mean of the two profiles and the dissipation uses their minimum,
\beeq
 \ell_k = \sqrt{\ell^{up}\, \ell^{dwn}} , \qquad \ell_\epsilon = \min \lk \ell^{up},\ \ell^{dwn} \rk
\eeeq
$\ell_k$ replaces $\ell$ in the diffusivity $\kappa_M$, and $\ell_\epsilon$ replaces $\ell$ in the dissipation $c_\epsilon e^{3/2}/\ell$.
The constant $c_\epsilon$ multiplies the dissipation here.
Writing the dissipation as $e^{3/2}/\ell_\epsilon$ with $\ell_\epsilon = \min(\ell^{up}, \ell^{dwn})/c_\epsilon$ is the same.
The surface length $\ell_0$ has two forms.
Without a roughness length it follows Charnock's relation
\beeq
 \ell_0 = \kappa_{vk}\, \beta\, |\v \tau| / g
\eeeq
with the kinematic surface stress $\v \tau$, i.e.\ the stress divided by $\rho_0$,
the von K\'arm\'an constant $\kappa_{vk} = 0.4$ (attribute \verb|mxl_kappa|)
and the Charnock constant $\beta$ (attribute \verb|mxl_beta|).
The default $\beta = 2\cdot 10^5$ is the value of NEMO. It gives $\ell_0 = 0.8\,$m at a stress of $0.1\,$Pa.
The value $\beta = 2\cdot 10^4$ gives $\ell_0 = 0.08\,$m at the same stress.
The modulus $|\v \tau|$ is formed pointwise from \verb|taux| and \verb|tauy|,
or taken from \verb|tau_modulus| if the configuration provides it.
With a roughness length $z_0 > 0$ (attribute \verb|z0_surf|) the surface length is $\ell_0 = \kappa_{vk}\, z_0$ instead.
The coupled configurations set $z_0 = 1.6\, H_s$ with the significant wave height $H_s$ from the wave model.
In both forms $\ell_0 \ge \ell_{min0}$.
One halo exchange of $\ell_0$ is made per time step.
A configuration whose stress carries valid halo values can avoid it by setting \verb|tau_halo_ok|.

The discretisation follows pyOM2: backward Euler with the dissipation
linearised in the diagonal ($c_\epsilon \sqrt{e^n}/\ell$ acting on
$e^{n+1}$), a surface $W$ cell of half thickness $\Delta z/2$, and the
surface flux boundary condition added to that cell,
\beq
 F_e(0) = \left\{
 \begin{array}{ll}
   \Phi_{oc} & \mbox{wave dissipation flux from the wave model} \\
   (|\v \tau|)^{3/2} & \mbox{classical parameterisation}
 \end{array} \right.
\eeq
selected by \verb|tke_surf_from_waves| (default on, falling back to the
classical form if no wave flux is present; $\v \tau$ kinematic; the first
form closes the energy cycle between the wave and ocean components, the
second is the pyOM2 default and serves as the control experiment).
The classical form carries the factor one of pyOM2.
Craig and Banner (1994) give $\alpha_{cb} |\v \tau|^{3/2}$ with $\alpha_{cb} \approx 100$.
pe\_simple has no attribute for this factor.
A configuration that wants it overloads \verb|_surface_forcing|,
as \verb|setups/coupled_common/exp_ocean.py| does with \verb|EXP_ALPHA_CB|. No
diffusive TKE, tracer or momentum flux passes through the surface face
itself: the surface forcing enters only as the explicit flux above, which
also removes a spurious shear production against the zero-velocity halo
above the surface (a phantom no-slip lid).
For wave-driven injection a surface roughness length $z_0$ can be set
(attribute \verb|z0_surf|, default $0$, e.g.\
$1.6\,H_s$ from the wave model in the sense of Mellor \& Blumberg 2004).
With \verb|mxl_choice| $=1$ it only raises the limit of the value at the surface face.
The downward sweep runs before that limit is applied, so the mixing length below the surface face does not change.
With \verb|mxl_choice| $=2$ it sets the surface length $\ell_0$ described above.
If the
surface cell is drained negative by buoyancy fluxes, the missing energy is
added back and booked in \verb|tke_surf_corr| as in pyOM2.
There is no other floor on $e$, as in pyOM2.
A floor $e \ge e_{min}$ is applied only if \verb|tke_min| $> 0$ is set.
The initial value of $e$ in all wet cells is \verb|tke_init| $= 10^{-12}\,$m$^2$/s$^2$, as in the pyOM2 configurations.

Vertical friction is
computed implicitly with $\kappa_M$ averaged to the $U$/$V$ points as in
pyOM2 ($\overline{\kappa_M}^i$).
A momentum flux exists only across a face between two wet $U$ (or $V$) cells, as in pyOM2.
TKE is likewise exchanged only between two wet $W$ levels.
The shear production is diagnosed
from the same discrete operator with the updated velocities,
\beq
 K_{diss} = \Bar{\, \overline{\kappa_M}^i \lk \pz{u^{n+1}} \rk^2 \,}^i
          + \Bar{\, \overline{\kappa_M}^j \lk \pz{v^{n+1}} \rk^2 \,}^j
\eeq
The backward-Euler friction step is $\v u^{n+1} - \v u^n = \Delta t\, \pz{} \lk \overline{\kappa_M}\, \pz{\v u^{n+1}} \rk$.
Multiplying it by $\v u^{n+1}$ and integrating over the column gives the kinetic energy it removes,
\beeq
 \int_{-h}^0 \frac{|\v u^n|^2 - |\v u^{n+1}|^2}{2 \Delta t}\, dz
 = \int_{-h}^0 \overline{\kappa_M} \lk \pz{\v u^{n+1}} \rk^2 dz
 + \int_{-h}^0 \frac{|\v u^{n+1} - \v u^n|^2}{2 \Delta t}\, dz
\eeeq
The first term on the right is $K_{diss}$ and is the forcing of the TKE equation.
The second term is the numerical dissipation of the backward-Euler step.
It is of order $\Delta t$ and is not put into the TKE equation.
The kinetic energy removed from the mean flow therefore reappears in the TKE forcing
up to this term. $P_{diss} = \kappa_H N^2$ is the
exact transfer of the implicit buoyancy mixing, evaluated with the
\emph{updated} buoyancy (the tracer solve precedes the TKE step; using the
pre-mixing $N^2$ over-drains the TKE whenever $\kappa_H \Delta t/\Delta z^2$
is not small).
At the surface face $P_{diss}$ is the surface buoyancy flux instead, as in pyOM2,
\beeq
 P_{diss}(0) = - \frac{g}{\rho_0} \lk \frac{\p \rho}{\p T} F_T + \frac{\p \rho}{\p S} F_S \rk
\eeeq
Here $F_T$ and $F_S$ are the surface fluxes of temperature and salt into the ocean
(K\,m/s and psu\,m/s, positive into the ocean).
The configuration provides them as the attributes \verb|forc_temp_surface| and \verb|forc_salt_surface|.
Surface cooling is thus a source of TKE.
Without these attributes the term is zero.
For \verb|eq_of_state| $=0$ the temperature is the buoyancy $b$ and $\p \rho / \p b = -\rho_0/g$ is used.
pyOM2 uses zero there and has no surface buoyancy source of TKE for this equation of state.
The surface $W$ cell has the thickness $\Delta z/2$.
The TKE that enters it per unit area and time is therefore
$F_e(0) + (\Delta z/2)\, J_s$ plus the correction \verb|tke_surf_corr| described above.
Here $J_s = - P_{diss}(0)$ is the surface buoyancy loss.
In stand-alone column
tests the discrete TKE budget (surface flux $+$ shear production $=$
dissipation $+$ buoyancy destruction $+$ storage) closes to machine
precision.

When the model velocity is the Lagrangian mean $\v u_L$ of the coupled
wave--ocean formulation, the attribute \verb|mix_eulerian| (default off,
requires the Stokes velocity to be present) lets the friction operator act
on the Eulerian mean velocity $\bar{\v u} = \v u_L - \v u_S$: the
down-gradient closure parameterises Eulerian fluctuations, so mixing
$\v u_L$ would spuriously diffuse the Stokes shear. The Stokes velocity is
subtracted before the implicit solve and added back afterwards, and the
shear-production booking becomes flux times \emph{Lagrangian} shear, split
into the non-negative Eulerian part $K_{diss}^E = \kappa_M (\pz{\bar{\v u}})^2$
and the sign-indefinite Stokes (Langmuir) part
$K_{diss}^L = \kappa_M \pz{\bar{\v u}} \cdot \pz{\v u_S}$, kept separately
in the diagnostics with $K_{diss} = K_{diss}^E + K_{diss}^L$ feeding the TKE
equation. Without \verb|mix_eulerian| this reduces bit-identically to the
plain booking above.

The calling sequence per time step is: diffusivities from the current
$e$ and $N^2$ at the start of the step (so buoyancy mixing and friction
use the same coefficients the transfers are paired with), vertical
friction after the pressure correction, and the TKE step at the end of
the time step.
The density is evaluated once per time step, as in pyOM2.
This happens after the vertical mixing of the tracers and before the TKE step.
$N^2$ and $P_{diss}$ of the TKE step use this updated density.
The hydrostatic pressure of the next time step is computed from the same density.


\section*{Model parameters and options}

All parameters and options are attributes of the model instance, set either
through the parameter dictionary passed at construction or by assigning the
attribute before \verb|prepare_for_run|; options with a default may be left
unset. The tables list every attribute the model reads.

\subsection*{Required parameters (no default)}

\begin{longtable}{@{}llp{8.4cm}@{}}
\hline
Attribute & Unit & Meaning \\
\hline
\endhead
\verb|nx|, \verb|ny|, \verb|nz| & --- & global grid size (cells, without halo
  points) \\
\verb|dx|, \verb|dy|, \verb|dz| & m & grid spacing; $\Delta z$ is a uniform
  scalar \\
\verb|dt| & s & time step \\
\verb|eq_of_state| & --- & equation of state: $0$ (temperature carries the
  buoyancy, no salinity), $1$ (linear $\rho(S,T)$, salinity prognostic),
  $5$ (GSW/TEOS-10, salinity prognostic, the pressure is the depth of the
  cell centre below the surface in dbar), $100$ (ideal-gas atmosphere) \\
\verb|implicit_free_surface| & --- & implicit free surface if true \\
\verb|explicit_free_surface| & --- & split-explicit free surface if true;
  both false: rigid lid \\
\verb|periodic_in_x|, \verb|periodic_in_y| & --- & lateral boundaries
  periodic (true) or closed (false) \\
\verb|epsilon| & --- & convergence threshold of the iterative surface-pressure
  solver (rigid lid / implicit free surface) \\
\verb|max_iterations| & --- & iteration cap of that solver \\
\verb|sub_cycling| & --- & number of barotropic sub-cycles $M$ (split-explicit
  scheme only) \\
\hline
\end{longtable}

\noindent
The constructor additionally takes \verb|n_pes_x|, \verb|n_pes_y| (processor
grid, default $1 \times 1$), \verb|halo_points| (default 1) and
\verb|comm_in| (default: MPI world communicator). The constants
\verb|grav| $=9.81$ and \verb|rho_0| $=1024$ are set internally (note: they
overwrite any value passed in the parameter dictionary). Coriolis parameter,
topography (\verb|kbot|) and initial conditions are set through the template
methods \verb|set_coriolis|, \verb|set_topography| and
\verb|set_initial_conditions|, overloaded in the configuration.

\subsection*{Forcing hooks and coupling fields}

\begin{longtable}{@{}lp{10.6cm}@{}}
\hline
Attribute / method & Meaning \\
\hline
\endhead
\verb|set_forcing(tau)| & optional user method; adds forcing to the momentum
  tendencies \verb|du[tau]|, \verb|dv[tau]| (enters the AB interpolation).
  Surface fluxes of heat and salt are added to \verb|dtemp[tau]| and
  \verb|dsalt[tau]| in the same way. Both schemes keep these tracer
  increments \\
\verb|set_forcing_direct()| & optional user method; direct momentum increment
  bypassing the AB weights, applied before the pressure solve. Used for the
  Stokes tendency $\Delta \v u_S$ of the Lagrangian-mean ($\v u_L$)
  formulation, which must be a plain forward increment.
  It is called only in the rigid-lid and implicit free-surface scheme.
  The split-explicit scheme does not call it \\
\verb|forc_temp_surface| & surface flux of temperature into the ocean
  (K\,m/s). It gives the surface buoyancy source $P_{diss}(0)$ of the TKE
  equation. It does not force the tracer. That is done in \verb|set_forcing| \\
\verb|forc_salt_surface| & the same for salt (psu\,m/s) \\
\verb|tau_modulus| & optional modulus of the kinematic surface stress at the
  $T$ points, used by \verb|mxl_choice| $=2$ instead of the pointwise modulus
  of \verb|taux|, \verb|tauy| \\
\verb|taux|, \verb|tauy| & kinematic wind stress; also the fallback surface
  TKE forcing $|\v \tau|^{3/2}$ \\
\verb|phi_oc| & wave-breaking energy flux from the wave model; surface TKE
  flux when \verb|tke_surf_from_waves| is on \\
\verb|u_S|, \verb|v_S| & Stokes velocity from the wave model; required by
  \verb|mix_eulerian| \\
\hline
\end{longtable}

\subsection*{Options with defaults}

\begin{longtable}{@{}llp{5.9cm}@{}}
\hline
Attribute & Default & Meaning \\
\hline
\endhead
\verb|superbee_advection| & unset & if set true, tracer advection uses the
  Superbee flux limiter instead of the second-order fluxes (momentum
  advection is always second order). Note: setting it to false raises an
  error; leave it unset for the second-order scheme \\
\verb|Khbi| & unset & biharmonic horizontal tracer diffusion coefficient
  (m$^4$/s); active if present, forward Euler \\
\verb|Ahbi| & unset & biharmonic horizontal friction coefficient (m$^4$/s);
  active if present, forward Euler. In the rigid-lid and implicit
  free-surface scheme it is applied to $\v u$ before the AB update. In the
  split-explicit scheme it is a tendency in the forcing $\v R^n$ of the
  external mode \\
\verb|vmix_scheme| & \verb|'explicit'| & vertical mixing/friction hierarchy:
  \verb|'explicit'| (forward Euler, coefficient checked against the
  diffusive stability limit) or \verb|'implicit'| (backward Euler,
  unconditionally stable). If unset, the legacy flag
  \verb|enable_implicit_vmix| decides \\
\verb|kappa_T| & $0$ & background vertical diffusivity (m$^2$/s) \\
\verb|A_v| & $0$ & background vertical viscosity (m$^2$/s) \\
\verb|enable_conv_adjust| & on iff implicit & convective adjustment:
  $\kappa = \kappa_{conv}$ where $N^2 < 0$ (implicit scheme only) \\
\verb|kappa_conv| & $1$ & convective-adjustment diffusivity (m$^2$/s) \\
\verb|enable_tke| & off & TKE closure; overrides the constant coefficients
  for both buoyancy and momentum (implicit scheme only) \\
\verb|enable_bottom_friction| & off & linear bottom friction on the lowest
  wet cell, rate \verb|r_bot| \\
\verb|r_bot| & $0$ & rate of the linear bottom friction (1/s) \\
\verb|enable_quadratic_bottom_friction| & off & quadratic bottom friction on
  the lowest wet cell, rate $C_d |\v u_b|/\Delta z$ \\
\verb|r_quad_bot| & $0$ & drag coefficient $C_d$ of the quadratic bottom
  friction \\
\verb|fb_theta| & $0.14$ & dissipation parameter $\theta$ of the
  split-explicit sub cycle \\
\verb|free_surface_tracer_flux| & off & tracer flux through the surface face.
  Off means zero flux, as in pyOM2. On means $w\, b$ with the surface-cell
  value, so a uniform tracer stays uniform \\
\hline
\end{longtable}

\subsection*{TKE closure constants (read at \texttt{init\_tke})}

\begin{longtable}{@{}llp{8.2cm}@{}}
\hline
Attribute & Default & Meaning \\
\hline
\endhead
\verb|c_k| & $0.1$ & diffusivity constant, $\kappa_M = c_k \ell \sqrt{e}$ \\
\verb|c_eps| & $0.7$ & dissipation constant, $c_\epsilon e^{3/2}/\ell$ \\
\verb|alpha_tke| & $1$ & vertical TKE diffusion $\alpha \kappa_M$; the pyOM2
  acc configurations use $30$ \\
\verb|mxl_min| & $10^{-12}$ & mixing-length floor (m); acc configurations use
  $10^{-8}$ \\
\verb|kappaM_min| & $0$ & viscosity floor (m$^2$/s), as in pyOM2. A background
  viscosity is set here \\
\verb|kappaM_max| & $100$ & viscosity cap (m$^2$/s) \\
\verb|kappaH_min| & $0$ & diffusivity floor (m$^2$/s). A background
  diffusivity is set here \\
\verb|tke_min| & $0$ & TKE floor (m$^2$/s$^2$), applied only if positive \\
\verb|tke_init| & $10^{-12}$ & initial TKE in wet cells (m$^2$/s$^2$) \\
\verb|mxl_choice| & $1$ & limits of the mixing length. $1$ is pyOM2,
  $2$ is NEMO 4.2 with a separate length for the dissipation \\
\verb|mxl_kappa| & $0.4$ & von Karman constant of the surface length
  (\verb|mxl_choice| $=2$) \\
\verb|mxl_beta| & $2\cdot 10^5$ & Charnock constant $\beta$ of the surface
  length (\verb|mxl_choice| $=2$) \\
\verb|mxl_min0| & $0.04$ & minimum surface and bottom length (m)
  (\verb|mxl_choice| $=2$) \\
\verb|tau_halo_ok| & off & set when \verb|taux|, \verb|tauy| carry valid halo
  values. Saves the halo exchange of the surface length
  (\verb|mxl_choice| $=2$) \\
\verb|z0_surf| & $0$ & surface roughness length (m), e.g.\ $1.6\,H_s$.
  Limits the value at the surface face (\verb|mxl_choice| $=1$) or sets the
  surface length (\verb|mxl_choice| $=2$) \\
\verb|tke_surf_from_waves| & on & surface TKE flux from the wave dissipation
  $\Phi_{oc}$; falls back to $|\v \tau|^{3/2}$ if no wave flux is present \\
\verb|mix_eulerian| & off & vertical friction acts on the Eulerian mean
  velocity $\v u_L - \v u_S$ ($\v u_L$ formulation), with the
  $K_{diss}^E$/$K_{diss}^L$ booking split \\
\hline
\end{longtable}

\subsection*{Restart state}

The model has no restart routine of its own. A configuration that continues a
run must save the following fields and set them again before the next step.
A restart is then bit-identical to the uninterrupted run.
This is tested by \verb|pyOM3/test/pe/restart.py|, which holds the reference list.
\begin{itemize}
\item Always the fields \verb|u|, \verb|v|, \verb|w|, \verb|temp|, \verb|salt|, \verb|rho|, \verb|p_h| and \verb|p_s|.
\item Always the three time levels of the tendencies \verb|du|, \verb|dv|, \verb|dtemp| and \verb|dsalt|.
\item The time-level pointers \verb|tau|, \verb|taum1|, \verb|taum2| and the AB coefficients of the driving loop.
The method \verb|loop| starts every call with a forward Euler step and the pointers $2,1,0$.
A restart must therefore drive the time steps itself.
\item With the TKE closure \verb|tke| and \verb|K_diss_v|.
The diffusivities of a step are built from $K_{diss}$ of the previous step through the Richardson number.
Without \verb|K_diss_v| the first step after the restart differs.
\item With the split-explicit scheme \verb|du_cor|, \verb|dv_cor|, \verb|U_bt| and \verb|V_bt|.
\item With \verb|mxl_choice| $=2$ \verb|taux| and \verb|tauy| (or \verb|tau_modulus|).
The closure uses them at the start of the step, before \verb|set_forcing| sets the new stress.
\end{itemize}

\subsection*{Snapshot diagnostics}

Two snapshot classes are available; either replaces the model class as the
base of the configuration and overrides \verb|prepare_for_run| and
\verb|diagnose|. Both read the attribute \verb|snapint| (snapshot interval in
seconds, required) and take the \verb|prepare_for_run| arguments
\verb|snap_file_name| and \verb|show_each_time_step| (default on: print every
time step; off: print at output steps only). When the iterative
surface-pressure solver is active, both print its per-step iteration count
and accumulate the total count between outputs.

\emph{snap\_nc4} writes a single netCDF file (default name
\verb|'snap.cdf'|), created with the T- and U-grid coordinate axes and the
variables $u$, $v$, $w$, temp, $p_h$, $\rho$ and $p_s$; salinity is included
when the equation of state carries it (\verb|eq_of_state| $>0$ and
$\ne 100$). At every output step each PE in turn writes its halo-free
sub-block into the global arrays of the shared file (serialised with
barriers). For the rigid-lid/implicit path the current time level of the
surface pressure is written, for the split-explicit path the plain 2D field.
This class works with the MPI domain decomposition and with
single-process JAX, but not in a multihost JAX environment, since all
processes must be able to open the same file.

\emph{snap\_raw} writes one \verb|.npy| file per variable, per PE and per
output step into a directory tree
\verb|<name>/pe=<rank>/n=<step>/| (default name \verb|'raw_snap'|; an
existing directory of that name is removed at start-up), with the
coordinate axes and the processor-grid layout
(\verb|decomposition.npy|) stored once at the top level. Each process only
touches its own subdirectory, and JAX sharded arrays are written through
their local shard, so this class also works in a multihost JAX
environment. The variable set is $u$, $v$, $w$, temp, salt, $p_h$ and
$p_s$ (salt always, $\rho$ not included). The fragmented output is meant
to be reassembled afterwards using the stored decomposition, e.g.\ into
zarr.


\end{document}


